\documentclass{article}
\usepackage{graphicx} % Required for inserting images

\title{Gravitational wave production using Simulation Based Inference}
%\author{Murman Gurgenidze}
\date{October 2026}

\begin{document}

\maketitle

\section{Introduction}
Magnetohydrodynamic (MHD) turbulence in the early Universe, whether stirred by a
first-order phase transition or by primordial magnetic fields, sources a
stochastic gravitational wave (GW) background that can be potentially
observable by current and future GW experiments like LISA and Pulsar
Timing Arrays.
Its signal is set by the unequal-time correlator (UETC) of the
anisotropic stress, $E_\Pi(k,\tau_1,\tau_2)$, which encodes all the information
required for a robust predicition of the GW spectrum.
The UETC depends on the non-linear dynamics of the source of GWs and, therefore,
it requires to conduct high-resolution numerical simulations of the sourcing process
to make an accurate estimate.
Furthermore, computing the UETC is computationally costly and most simulations
usually just compute the resulting GW background by solving the GW equation numerically,
which is more efficient than computing UETCs of the sourcing fields
(e.g., velocity or magnetic fields in sources of GWs due to the development
of MHD turbulence in the early Universe).
On the other hand, knowledge of the UETC would allow us to develop theoretical
models that can be applied for a range of parameters that characterize the source
of GWs, avoiding running expensive simulations for each input parameter
of interest.

The objective of the propject is to use simulation-based inference to use
the results of numerical simulations that provide the GW spectrum from MHD
turbulent processes to train neural networks that allow to develop templates
for the GW spectrum and for the UETC of the source.
While the former is already enough to make predictions and parameter reconstruction
given the observation, the latter modeling can provide us with a more thorough
theoretical understanding of the sourcing process, making the results more
trustable, falsifiable, and more directly applicable to a broader range
of the input parameters.

Due to the humongous number of sources that we expect at the LISA band, the
LISA Global Fit combining all the different possible sources, both transient and
stochastic, as well as reconstruction of the different sources of noise, makes the
problem of reconstructing a stochastic GW background originating from early Universe
sources extremelly challenging.
Therefore, accurate but fast reconstruction of the GW background from the parameters
that
characterize the early Universe process leading to the GWs as, for example,
a first-order phase transition around the electroweak scale, will become indispensable
for a potential study of these signals from an observed GW background.
Hence, these results could be key for studying different beyond the Standard Model
physics that lead to a first-order phase transition, potentially providing
an explanation to other fundamental open problems like the origin of baryons 
over antibaryons.y

We propose an MHD-only program that makes UETCs fast while building in the physics.
We will learn the space--time spectrum $\tilde{E}_\Pi(k,\omega;T)$ of kinetic,
magnetic and cross stresses directly. Positivity of the GW spectrum then holds by
construction, and the signal follows from this spectrum on the light cone,
$\omega=k$, where Alfv\'enic resonances and nonlinear decorrelation separate
cleanly. We will
(i) build streaming UETC diagnostics for MHD lattice codes;
(ii) learn reduced stochastic dynamics that extrapolate to late times and high Reynolds numbers;
(iii) quantify non-Gaussian corrections from intermittent magnetic fields; and
(iv) build differentiable emulators in a compact, source-agnostic parameter space,
including GW polarization from helical turbulence.

These emulators make simulation-based inference (SBI) possible. Our SBI pipeline
will treat turbulence efficiency and UETC model uncertainty as nuisance
parameters, so posteriors honestly include theory error. It will combine our
turbulence emulator with existing sound-wave models, handle realistic foregrounds,
and allow joint inference with intergalactic magnetic-field constraints.

The outcome will be validated MHD UETC models, a public correlator dataset,
millisecond emulators, and an inference pipeline that links GW observations to
early-universe turbulence and primordial magnetism with quantified uncertainty.

\end{document}